\section{True real degree tests on the original closed integral curves}
\label{sec:native-closed-curve-degrees}

\begin{theorem}[Actual original closed-curve degree functionals]
\label{thm:native-closed-curve-degrees}
\lean{EnvelopeProperty.NativeClosedCurveDegrees.ClosedIntegralCurve, EnvelopeProperty.NativeClosedCurveDegrees.ClosedIntegralCurve.structureMorphism, EnvelopeProperty.NativeClosedCurveDegrees.ClosedIntegralCurve.projectiveFactorization, EnvelopeProperty.NativeClosedCurveDegrees.ClosedIntegralCurve.picardDegree, EnvelopeProperty.NativeClosedCurveDegrees.ClosedIntegralCurve.picardDegree_lineClass, EnvelopeProperty.NativeClosedCurveDegrees.ClosedIntegralCurve.additiveDegree, EnvelopeProperty.NativeClosedCurveDegrees.ClosedIntegralCurve.realDegree, EnvelopeProperty.NativeClosedCurveDegrees.ClosedIntegralCurve.realDegree_tmul, EnvelopeProperty.NativeClosedCurveDegrees.ClosedIntegralCurve.realDegree_lineClass}
\leanok
\uses{thm:native-real-picard, thm:fc-original-picard-curve-degree}
\source{bj22-trivval:sec:tc}
For an original native-projective $X$ over the original field,
every actual closed integral one-dimensional curve in that same
original $X$ gives a genuine real linear degree functional on
the real scalar extension of its original native Picard group.
Its value on every actual original invertible line is precisely
the true whole Euler degree of that same line's actual pullback
to that same original curve. The original curve projection's
projectivity is derived from the original closed immersion and
original projectivity of $X$; no curve-degree or comparison
certificate is supplied.
\end{theorem}

\begin{proof}
The curve test consists exactly of an original integral scheme,
its actual closed immersion into the original $X$, and its true
topological dimension one. Compose that same immersion with the
original coefficient map of $X$. The actual finite polynomial
factorization of the original projective $X$, composed with
this same original closed immersion, derives the true projective
factorization of this original curve. The actual original
Picard pullback and full Euler curve-degree homomorphism therefore
give an integer degree homomorphism on the same native original
Picard group. Its value on an original line is the true whole
Euler degree of the actual original pullback, by the earlier
native line-class comparison.

Pass to the same associated additive Picard group, compose this
true integer homomorphism with the original integer inclusion in
the reals, and use the native tensor scalar-extension universal
property. The resulting real linear functional evaluates
$r\otimes a$ as $r$ times the same true original curve degree.
In particular $1\otimes[L]$ evaluates as the true original full
Euler degree of the unchanged original line pullback. These are
exact original closed curves, as in the source definition of
numerical equivalence. No equivalence with tests by arbitrary
curve morphisms, no finite-dimensionality, and no numerical
quotient or ample-cone theorem is asserted here; those require
their actual geometric producers.
\end{proof}
